\chapter*{Заключение}
\addcontentsline{toc}{chapter}{Заключение}

В настоящей диссертации разработаны методы информационного поиска для русскоязычных больших языковых моделей и предложена их интеграция в систему обеспечения безопасности мультимодальных моделей. Сквозная линия работы --- от поискового стека к безопасному индексу и далее к двухэтапному конвейеру безопасности --- объединяет извлечение достоверной информации с механизмами модерации контента в единую систему.

\textbf{Основные результаты.} В ходе работы получены следующие результаты.
\begin{enumerate}
    \item Систематизированы современные подходы к построению текстовых и мультимодальных эмбеддингов, методам переранжирования, парадигмам генерации с извлечением информации, а также к системам обеспечения безопасности и уязвимостям мультимодальных моделей (Глава~\ref{ch:lit_review}).
    \item Представлен фреймворк GigaEmbeddings: иерархический трёхэтапный конвейер инструктивного дообучения декодерной LLM GigaChat-3B с двунаправленным вниманием и латентным пулингом. Модель занимает первое место на русскоязычном бенчмарке MTEB(rus, v1.1) со средним баллом 74.16, превосходя сильные базовые модели с большим числом параметров (Глава~\ref{ch:gigaembeddings}).
    \item Предложен кросс-энкодер реранкер на основе компактных декодерных моделей GigaChat-0.6B и GigaChat-3B, адаптирующий декодер в реранкер заменой каузального внимания на двунаправленное, агрегацией mean pooling и выделенной классификационной головой. На трёх русскоязычных IR-бенчмарках реранкер GigaChat-3B достигает наилучшего качества среди открытых моделей сопоставимого размера (Глава~\ref{ch:reranking}).
    \item Построен двухэтапный поисковый стек, объединяющий лексический (BM25), плотный, гибридный и двухэтапный поиск, и обоснованы преимущества RAG в части снижения галлюцинаций и устранения необходимости переобучения при обновлении данных (Глава~\ref{ch:rag}).
    \item Разработан двухэтапный конвейер обеспечения безопасности мультимодальных моделей, включающий потоковый классификатор обменов на основе MoE-модели GigaChat3.1-10B-A1.8B и безопасную модель, дообученную методом конституционного ИИ и обучения с подкреплением с гибридной наградой за безопасность. Поисковый стек интегрирован в конвейер в виде инструмента безопасного поиска по курируемому индексу (Глава~\ref{ch:safety}).
\end{enumerate}

\textbf{Ключевой методологический вывод.} Эмпирически (Глава~\ref{ch:safety}) показано, что присвоение метки безопасности обмену в целом, а не изолированному запросу, необходимо для надёжной фильтрации: при переходе к разметке по фактической безопасности порождённого ответа качество классификатора запросов падает до уровня, близкого к случайному ($\mathrm{AUC} = 0.60$), тогда как потоковый классификатор ответа сохраняет существенно более высокую различающую способность ($\mathrm{AUC} = 0.83$). Этот результат обосновывает выбор потоковой классификации обмена в качестве первого этапа конвейера.

\textbf{Ограничения.} Эмпирически валидированными результатами автора являются модель эмбеддингов (Глава~\ref{ch:gigaembeddings}), реранкер (Глава~\ref{ch:reranking}, со-руководство), сравнение схем классификатора безопасности и конструкция награды (Глава~\ref{ch:safety}). Конвейер безопасности является мультимодальным по построению: потоковый классификатор оценивает изображения, вызовы инструментов и текст, а защищаемая модель обрабатывает изображения, использует режим рассуждения и обращается к инструменту безопасного поиска для запросов любой модальности. Вместе с тем сквозная количественная оценка конвейера в целом --- доля блокировок и чрезмерных отказов, доля высокорисковых уязвимостей, бюджет задержки защитного слоя --- на русскоязычных и мультимодальных бенчмарках безопасности выходит за рамки настоящей работы и составляет предмет дальнейших исследований.

\textbf{Направления дальнейших исследований.} Перспективными представляются: сквозная экспериментальная валидация мультимодального конвейера безопасности на русскоязычных и мультимодальных бенчмарках безопасности (доля блокировок, доля чрезмерных отказов, доля высокорисковых уязвимостей); измерение сквозного бюджета задержки защитного слоя; а также распространение предложенной адаптации декодерных моделей (двунаправленное внимание, mean pooling, классификационная голова) на смежные задачи обработки естественного языка.
